% J2.5: propagación de incertidumbre por Monte Carlo en el edificio de tres plantas.
\documentclass[tikz,border=4pt]{standalone}
\usepackage{resonancia}
\begin{document}
\begin{tikzpicture}
  \node[cajaAzul, minimum width=3.2cm, minimum height=1.5cm] (e) at (0,0) {\textbf{entradas inciertas}\\[-1pt]{\footnotesize $k_i\sim N(\mu_i,\ 5\text{\ \%})$}\\[-1pt]{\footnotesize $m_i\sim N(\mu,\ 2\text{\ \%})$}};
  \node[cajaVerde, minimum width=3.4cm, minimum height=1.5cm] (m) at (5.0,0) {\textbf{modelo}\\[-1pt]{\footnotesize $\mathbf K\boldsymbol\varphi=\omega^2\mathbf M\boldsymbol\varphi$}\\[-1pt]{\footnotesize (una vez por muestra)}};
  \node[cajaNaranja, minimum width=3.2cm, minimum height=1.5cm] (s) at (10.0,0) {\textbf{salidas}\\[-1pt]{\footnotesize $f_1,f_2,f_3$: media,}\\[-1pt]{\footnotesize desviación, intervalo 95 \%}};
  \draw[flecha] (e) -- node[etiqueta, above] {20\,000 muestras} (m);
  \draw[flecha] (m) -- (s);
  \node[etiqueta, color=rGris, anchor=north] at (5.0,-1.1) {contraste: propagación lineal con derivadas, $\sigma_f^2=\sum_j (\partial f/\partial\theta_j)^2\sigma_j^2$};
\end{tikzpicture}
\end{document}
